%%=============================================================================
%% Conclusie
%%=============================================================================

\chapter{Conclusie}%
\label{ch:conclusie}

% TODO: Trek een duidelijke conclusie, in de vorm van een antwoord op de
% onderzoeksvra(a)g(en). Wat was jouw bijdrage aan het onderzoeksdomein en
% hoe biedt dit meerwaarde aan het vakgebied/doelgroep? 
% Reflecteer kritisch over het resultaat. In Engelse teksten wordt deze sectie
% ``Discussion'' genoemd. Had je deze uitkomst verwacht? Zijn er zaken die nog
% niet duidelijk zijn?
% Heeft het onderzoek geleid tot nieuwe vragen die uitnodigen tot verder 
%onderzoek?

Deze bachelorproef onderzocht in welke mate bestaande AI-modellen inzetbaar zijn voor automatische productherkenning in supermarktbeelden. De focus lag hierbij op de technische haalbaarheid van een reproduceerbare pipeline voor objectdetectie en similarity-based productidentificatie binnen retailomgevingen met hoge objectdichtheid.

Uit de literatuurstudie bleek dat automatische productherkenning in supermarkten een uitdagend probleem blijft. Producten staan vaak dicht bij elkaar, overlappen elkaar regelmatig en verschillen soms slechts in kleine details zoals kleuraccenten, tekst of verpakkingsvarianten. Daarnaast zorgen variaties in belichting, camerahoek en beeldkwaliteit voor bijkomende moeilijkheden. Hoewel moderne computer vision-modellen sterke prestaties tonen in algemene detectietaken, blijkt herkenning op SKU-niveau in realistische winkelomgevingen nog steeds complex.

Om deze problematiek te onderzoeken werd een proof-of-concept ontwikkeld bestaande uit meerdere componenten. Voor objectdetectie werd gebruik gemaakt van YOLO11s, terwijl productidentificatie gebeurde via embeddings gebaseerd op CLIP, OCR en Florence-2. De embeddings werden gecombineerd en vervolgens doorzocht via FAISS similarity search om mogelijke productmatches terug te geven.

De experimentele evaluatie toont aan dat fine-tuning een zeer grote impact heeft op detectieprestaties binnen retailomgevingen. Het generieke YOLO11s-model presteerde onvoldoende op supermarktbeelden met hoge objectdichtheid, maar na fine-tuning op de SKU-110K dataset steeg het aantal detecties aanzienlijk. Dit bevestigt dat modelspecialisatie noodzakelijk is voor toepassingen binnen complexe supermarktcontexten.

Daarnaast werd onderzocht welke combinatie van embeddings het meest geschikt is voor similarity-based productherkenning. Uit de experimenten bleek dat de combinatie van CLIP en OCR de meest consistente resultaten opleverde. Vooral tekstinformatie afkomstig van productverpakkingen bleek een belangrijke bijdrage te leveren aan correcte productmatching. De visuele encoder van Florence-2 bood daarentegen weinig meerwaarde binnen deze pipeline en werd uiteindelijk niet opgenomen in de finale configuratie.

Hoewel de top-5 resultaten vaak relevante producten bevatten, bleef de stabiliteit van top-1 voorspellingen beperkt. Producten met sterk gelijkende verpakkingen, zoals varianten van hetzelfde merk, bleken moeilijk betrouwbaar van elkaar te onderscheiden. Kleine verschillen in kleur, tekst of layout waren vaak onvoldoende om consistente embeddings te genereren. De vooropgestelde prestatiedrempels werden bovendien niet gehaald voor fine-grained productidentificatie, wat aantoont dat SKU-level herkenning in realistische retailomgevingen nog steeds een uitdagend probleem blijft. Dit benadrukt dat multimodale embeddings en similarity search reeds bruikbare resultaten opleveren, maar momenteel nog onvoldoende robuust zijn voor volledig betrouwbare identificatie op productniveau.

De resultaten tonen ook aan dat de kwaliteit van eerdere stappen in de pipeline een directe invloed heeft op het eindresultaat. Onnauwkeurige bounding boxes leiden tot slechte crops, wat vervolgens minder betrouwbare embeddings veroorzaakt. Daarnaast hadden OCR-fouten en inconsistente datasetafbeeldingen een negatieve impact op similarity matching. De prestaties van het volledige systeem worden dus bepaald door de interactie tussen detectie, embeddings, OCR en datasetkwaliteit.

Ondanks deze beperkingen toont deze bachelorproef aan dat automatische assortimentsinventarisatie technisch haalbaar is als onderzoeksrichting. Moderne computer vision-modellen kunnen reeds een groot deel van de producten in complexe supermarktbeelden detecteren en gedeeltelijk identificeren. Vooral de combinatie van objectdetectie met multimodale embeddings biedt interessante mogelijkheden voor toekomstige toepassingen binnen retailanalyse, voorraadbeheer en monitoring van gezonde of duurzame producten.

Toch blijven er nog verschillende beperkingen die verder onderzoek vereisen. De gebruikte datasets bevatten niet uitsluitend Vlaamse supermarktbeelden, waardoor de generaliseerbaarheid naar specifieke winkelcontexten beperkt blijft. Daarnaast bleef de evaluatie van retrievalprestaties relatief beperkt tot similarity scores en top-k matches. Hierdoor blijft het moeilijk om de retrievalkwaliteit volledig objectief te vergelijken met andere systemen. Toekomstig onderzoek kan uitgebreidere retrievalmetrics gebruiken, zoals Recall@k, Precision@k of Mean Reciprocal Rank, om productidentificatie nauwkeuriger te evalueren.

Verder onderzoek kan zich richten op krachtigere multimodale modellen, verbeterde OCR-methoden en grotere gespecialiseerde retaildatasets. Ook realtime verwerking op consumentenhardware blijft een belangrijk aandachtspunt wanneer dergelijke systemen praktisch inzetbaar moeten worden. Daarnaast kan onderzocht worden in welke mate gespecialiseerde retaildatasets uit Vlaamse supermarktcontexten de prestaties van SKU-level herkenning verder kunnen verbeteren.

Samenvattend kan gesteld worden dat bestaande AI-modellen reeds sterke mogelijkheden bieden voor automatische productdetectie in supermarktbeelden, maar dat betrouwbare productherkenning op SKU-niveau nog aanzienlijke uitdagingen bevat. Deze bachelorproef toont echter aan dat een combinatie van objectdetectie, embeddings en multimodale informatie een veelbelovende basis vormt voor verdere ontwikkeling binnen automatische retailanalyse.